% Generated draft: requires mathematical and source review.
\documentclass[11pt]{article}
\usepackage[margin=25mm]{geometry}
\usepackage{fontspec}
\setmainfont{Latin Modern Roman}
\usepackage{amsmath,amssymb,mathtools,unicode-math}
\setmathfont{Latin Modern Math}
\usepackage{enumitem,needspace,xurl,hyperref}
\setlength{\parindent}{0pt}
\setlength{\parskip}{5pt}
\setlength{\emergencystretch}{3em}
\newcommand{\N}{\mathbb{N}}
\newcommand{\Z}{\mathbb{Z}}
\newcommand{\Q}{\mathbb{Q}}
\newcommand{\R}{\mathbb{R}}
\newcommand{\C}{\mathbb{C}}
\newcommand{\F}{\mathbb{F}}
\newcommand{\A}{\mathbb{A}}
\newcommand{\PP}{\mathbb{P}}
\newcommand{\E}{\mathbb{E}}
\newcommand{\eps}{\varepsilon}
\newcommand{\dd}{\,\mathrm{d}}
\begin{document}
\section*{UnsolvedMath ID: 30005678}
\textbf{Source problem code:} catalogue entry
\subsection{Definitions and mathematical statement}
In the classical setting, a Waldhausen category \(\mathcal C\) is a category with a zero object, a distinguished class of cofibrations, and a distinguished class of weak equivalences, satisfying the Waldhausen axioms, including existence and stability of pushouts along cofibrations and the gluing axiom for weak equivalences. The record does not specify whether classical categories or a higher-categorical version are intended.

For a classical Waldhausen category, the category \(S_n\mathcal C\) consists of diagrams \((A_{ij})_{0\leq i\leq j\leq n}\) with zero diagonal, such that, for every \(i\leq j\leq k\), the map \(A_{ij}\to A_{ik}\) is a cofibration and the square
\[
\begin{matrix}
A_{ij}&\longrightarrow&A_{ik}\\
\big\downarrow&&\big\downarrow\\
0&\longrightarrow&A_{jk}
\end{matrix}
\]
is a pushout. Morphisms are natural transformations; weak equivalences are objectwise weak equivalences. Restriction along order-preserving maps gives the simplicial \(S_\bullet\)-construction.

To describe the full \(2\)-Segal condition precisely in one possible framework, let \(X\colon\Delta^{\mathrm{op}}\to\mathrm{Spaces}\) be a simplicial space. For an integer \(n\geq2\) and a triangulation \(T\) of the polygon with ordered vertices \(0,\ldots,n\), let \(X(T)\) be the homotopy limit of the diagram assigning \(X_2\) to each triangle, \(X_1\) to each edge, and \(X_0\) to each vertex, with maps induced by face inclusions. Then \(X\) is fully \(2\)-Segal if, for every such \(n\) and \(T\), the restriction map
\[
\rho_T\colon X_n\longrightarrow X(T)
\]
is an equivalence.

\textbf{Extracted problem, with unresolved conventions.} Characterize the Waldhausen categories \(\mathcal C\) for which the intended realization \(X(\mathcal C)\) of \(S_\bullet\mathcal C\) is fully \(2\)-Segal, and find natural examples for which it is left \(2\)-Segal but not fully \(2\)-Segal. In the simplicial-space framework just described, failure of the full condition means that there exist \(n\geq3\) and a triangulation \(T\) for which \(\rho_T\) is not an equivalence.

The supplied record does not specify the realization \(X(\mathcal C)\), the equivalences or pullbacks used in the intended framework, or the definition and orientation convention for \emph{left} \(2\)-Segal. For example, taking nerves of the categories \(S_n\mathcal C\) and taking nerves of their weak-equivalence subcategories are distinct constructions. Thus the simplicial-space formulation above is conditional, and the left-only part cannot be made unambiguous from the supplied record alone. The word ``natural'' is also not given a formal criterion.

The supplied report excerpt motivates the question by contrasting categories with cofibrations and pullback data with Waldhausen categories, which have cofibrations and suitable pushouts but need not have an additional distinguished vertical class of morphisms or pullback data. This contextual comparison does not supply the missing conventions or establish a left \(2\)-Segal theorem for every intended Waldhausen category.
\subsection{Short English statement}
Determine which Waldhausen \(S_\bullet\)-constructions satisfy every polygonal \(2\)-Segal comparison, and seek natural examples satisfying only the intended left \(2\)-Segal condition. The source leaves the realization of the construction and the left \(2\)-Segal convention unspecified, so this is a qualified statement-only draft requiring mathematical review.
\subsection{Sources}
\begin{itemize}
\item Supplied local UnsolvedMath registry record, permanent ID \(30005678\), title ``Fully Two-Segal Waldhausen S-Constructions,'' exact source problem code \texttt{catalogue entry}, in the Oberwolfach Reports Open Problems collection. No individual proposer is supplied; the proposed year is recorded as 2023.
\item Supplied source citation: \emph{Homotopy Theory (2024)}, Oberwolfach Reports, DOI \texttt{10.4171/owr/2023/34}, pp.\ 1893--1952. Supplied source and literature-review link: \url{https://doi.org/10.4171/owr/2023/34}. The citation's title year 2024 and its DOI year component 2023 are retained as supplied.
\item Literature assessment dated exactly \texttt{2026-08-22}: the supplied curation labels the problem open and says that the extension from categories with cofibrations and pullbacks to general Waldhausen categories remains a natural question; the cited equivariant and scissors-congruence results are said not to settle that generalization. This is supplied dated curation, not an independently verified assessment. The record explicitly requests independent verification of time-sensitive status and cited claims.
\item The supplied registry labels the statement \texttt{corrected\_verified} and reports that it was reconstructed and checked against the official Oberwolfach report. These are source-record assertions; this draft does not independently verify them or claim to have visited the report.
\end{itemize}
\end{document}
